% !PRINTED 726
\weberpage{726}

\begin{center}
\S.\,199.

\gesperrt{The class number factor} $B$.
\end{center}

The class number factor $B$ must be calculated in quite another way. Here one has to set, if $\Theta$ is again a primitive root
of the equation \S.\,198, (3) ($\Theta^{\nu'} + 1 = 0$),
\weq{1}{t \equiv (-1)^\alpha\,5^\beta \pmod{m},}
\weq{2}{\chi_1(t) = \Theta^\beta}
and formula \S.\,197, (15) then gives:
\weq{3}{X_1 = -\frac{2}{m}\,(1, \Theta^{-1}, r)\,\overset{t}{\Sigma}\,\Theta^\beta\,\log\,\mathrm{tg}\,\frac{t\pi}{m}.}

Now for every odd $t$, if [with regard to (1)]
\[
r = e^{\frac{2\pi i}{m}}, \quad r^\nu = i, \quad r^{\nu t} = (-1)^\alpha i
\]
is set, we have
\weq{4}{\mathrm{tg}\,\frac{t\pi}{m} = (-1)^\alpha\,r^{t\nu}\,\frac{1 - r^t}{1 + r^t},}
and this, as we have already seen in \S.\,178, is a unit of the real field $H_m$. One obtains the conjugate values
of this unit with respect to this field by setting $t \equiv 5^\beta \pmod{m}$ and letting $\beta$ run through a complete system of residues with respect to the modulus $\mu$.
We accordingly set:
\weq{5}{\tau_\beta = \mathrm{tg}\,\frac{5^\beta\pi}{m}.}

By \S.\,198, (7) we have $5^{\beta+\nu'} \equiv 5^\beta + \mu \pmod{m}$, and from this
there follows, by a well-known trigonometric formula:
\weq{6}{\tau_{\beta+\nu'} = \frac{-1}{\tau_\beta}.}

If we therefore now set
\weq{7}{\lambda_\beta = \tfrac{1}{2}\,\log\,\tau_\beta^2,}
then by (6)
\weq{8}{\lambda_{\beta+\nu'} = -\lambda_\beta,}
and accordingly
\weq{9}{\Theta^{\beta+\nu'}\,\lambda_{\beta+\nu'} = \Theta^\beta\,\lambda_\beta;}
among the values of $\beta$ occurring in the sum (3)
no two are, as we have already seen in the preceding section,
congruent with respect to the modulus $\nu'$, and if we therefore by means of
%CONT>>
